%% =====================================================================
%%  Corrigés — Bilan 1 « Premiers contacts »
%% =====================================================================
\corrigesBilan{1}

\begin{corrige}{1}
\begin{reponses}[3]
\item \rep{is}
\item \rep{isn't}
\item \rep{are}
\item \rep{am}
\item \rep{is}
\item \rep{isn't}
\end{reponses}
\pourquoi{le sens impose la négation : Léa n'est pas américaine, Mrs.\ Lopez n'est pas mexicaine
(elle est américaine). \en{my English} est au singulier : \en{is}.}
\end{corrige}

\begin{corrige}{2}
\begin{reponses}[2]
\item \rep{Where are you from?}
\item \rep{How old are you?}
\item \rep{What's your name?}
\item \rep{Who's that?}
\item \rep{Yes, she is.}
\item \rep{No, they aren't.}
\end{reponses}
\pourquoi{acceptez aussi les formes pleines : \en{What is your name?}, \en{Who is that?}, et
\en{No, they're not.} pour la dernière réponse.}
\piege{\en{Yes, she's.} : une réponse courte affirmative n'est jamais contractée.}
\end{corrige}

\begin{corrige}{3}
\begin{reponses}[4]
\item \rep{his}
\item \rep{her}
\item \rep{their}
\item \rep{her}
\end{reponses}
\retiens{le possessif dépend du possesseur : Sam → \en{his} ; Mrs.\ Lopez et Léa → \en{her} ;
Sofia et David → \en{their}.}
\end{corrige}

\begin{corrige}{4}
\begin{reponses}[2]
\item \rep{Sam's cat}
\item \rep{Mrs. Lopez's daughter}
\item \rep{Emma's brother}
\item \rep{the boys' names}
\end{reponses}
\pourquoi{\en{boys} est un pluriel en \en{-s} : on ajoute seulement l'apostrophe, \en{the boys'
names}.}
\end{corrige}

\begin{corrige}{5}
\begin{reponses}[4]
\item \rep{fourteen}
\item \rep{forty}
\item \rep{nineteen}
\item \rep{seventy-two}
\end{reponses}
\piege{\en{fourty} : le \en{u} disparaît dans \en{forty} (mais pas dans \en{fourteen}).}
\end{corrige}

\begin{corrige}{6}
\begin{reponses}[3][french]
\item \rep{F}
\item \rep{V}
\item \rep{F}
\item \rep{V}
\item \rep{F}
\item \rep{V}
\end{reponses}
\pourquoi{c'est Léa qui est la voisine de Sam (\en{this is Sam, my neighbor}) ; Emma vient de
Londres ; Maya a vingt ans (\en{She's twenty}) ; elle dit \en{Welcome to my brother's party!}
Texte de la piste page~\pageref{trans-B1.01}.}
\end{corrige}

\begin{corrige}{7}
\begin{reponses}[3][french]
\item \rep{V}
\item \rep{F}
\item \rep{V}
\item \rep{V}
\item \rep{F}
\item \rep{F}
\end{reponses}
\pourquoi{Sam vient de Brooklyn (ses parents sont à Atlanta) ; Mango est un chat (\en{Our cat's
name is Mango}) ; Sam ne donne pas son numéro : \en{it's a secret!}}
\end{corrige}

\begin{corrige}{8}
Réponse modèle :
\begin{anglais}
\begin{tcolorbox}[otBase,colback=otGreenTint,colframe=otGreenTint,fontupper=\fontsize{10.4}{14}\selectfont]
Hi! My name is Léa Martin. I'm twenty-seven. I'm French, from Lyon. I speak French, English and a
little Italian. My parents and my brother are in Lyon.
\end{tcolorbox}
\end{anglais}
\retiens{comptez 0,5 point par information correcte (nom, âge, origine, langues, famille), dans
la limite de 2 points. Une information compte si la phrase est juste : \en{be} bien conjugué,
majuscules aux nationalités et aux langues.}
\end{corrige}

\begin{corrige}{9}
Réponses modèles :
\begin{reponses}[1]
\item \rep{Hi, Maya! I'm Léa.} \textit{ou} \rep{Hello! My name is Paul.}
\item \rep{Nice to meet you too. I'm from France.}
\item \rep{Yes, I am.}
\item \rep{No, it isn't. It's in France.}
\item \rep{Thank you!} \textit{ou} \rep{Thanks, Maya!}
\end{reponses}
\end{corrige}
